\section{Interpolation of diagonal form comparisons}\label{sec:interpolation}

The proof requires an interpolation principle for quadratic forms. Its
hypotheses involve four simple weights. In the later application at the
regularized minima, the target weight makes the squared norm of a centered
thermal-balance defect equal to the normalized entropy-production pairing
plus a centered-energy correction. This is the identity used in the
stationarity cancellation. We prove the principle here, including the
analytic property on which it depends.

Set
\begin{equation}\label{int:fb}
 f(x)=\frac{x}{\sinh x},\qquad b(x)=x\coth x,\qquad f(0)=b(0)=1.
\end{equation}
Both functions are positive and even on $\R$.  On Hilbert spaces
$\mathcal H_j$, $0\le j\le k$, fix coordinate representations in which all
weights below act by scalar multiplication.  A positive weight $X_j$
defines the possibly unbounded form
\[
 Q_{X,j}(Z)=\sum_{\alpha}X_j(\alpha)|Z_\alpha|^2;
\]
fixed positive coordinate measures can be absorbed into the coordinates.
Let $\mathcal T\subset\mathcal H_0$ be a linear test space, let
$L_i:\mathcal T\to\mathcal H_i$ be linear, and let $w_i>0$ for $1\le i\le k$.
We say that the family $X=(X_j)_{j=0}^k$ \emph{satisfies comparison} if
\begin{equation}\label{int:comparison}
 Q_{X,0}(Z)\ge\sum_{i=1}^k w_i Q_{X,i}(L_iZ)
 \qquad(Z\in\mathcal T).
\end{equation}
All sums on the right are interpreted as nonnegative coordinate sums.

\begin{theorem}[Four-weight interpolation]\label{thm:interpolation}
At every coordinate of every space, let $m>0$ and $x,y\in\R$ be given.
Suppose each $Z\in\mathcal T$ is supported on finitely many values of the
parameter tuple $(m,x,y)$, and $\mathcal T$ is invariant under scalar
multiplication by functions of this tuple.  If the four families
\begin{equation}\label{int:four-weights}
 mf(x)e^x,\quad mf(x)e^{-x},\quad mf(y)e^y,\quad mf(y)e^{-y}
\end{equation}
satisfy comparison, then so does
\begin{equation}\label{int:target-weight}
 F=m\frac{f(x)f(y)}{f(x+y)}.
\end{equation}
\end{theorem}

Finitely many parameter values do not mean finitely many nonzero
coordinates.  In particular, a centered finite matrix can have an infinite
identity tail, all at the same parameter value.  The hypothesis permits
this situation, which will occur in our application.

\subsection{Operations that preserve comparison}\label{int:operations}

The operations below belong to the parallel-sum and operator-mean
calculus of Anderson--Duffin and Kubo--Ando
\cite{AndersonDuffin1969,KuboAndo1980}. We give direct form proofs,
including the limits needed for unbounded weights.

Throughout this subsection, the auxiliary weights are positive functions
of the parameter tuple $(m,x,y)$.
Positive linear combinations preserve \eqref{int:comparison}.  So do
pointwise limits of such weights: on a source test there are
only finitely many parameter values, so its form converges; on the target
spaces Fatou's lemma gives the required inequality.

Comparisons also persist under the \emph{parallel sum}
\begin{equation}\label{int:parallel}
 X:Y=(X^{-1}+Y^{-1})^{-1}.
\end{equation}
Indeed, coordinatewise,
\[
 (X:Y)|z|^2=\min_{u+v=z}\bigl(X|u|^2+Y|v|^2\bigr),
 \qquad
 u=\frac{Y}{X+Y}z,\quad v=\frac{X}{X+Y}z.
\]
The source minimizers $U,V$ lie in $\mathcal T$.  Apply the comparisons for
$X$ and $Y$ to them, then use $L_iU+L_iV=L_iZ$ and the coordinatewise
minimum on each target space.  Target minimizers need not belong to any
specified test space.

Consequently, if $X,Y$ satisfy comparison, so do their logarithmic mean
$\ell(X,Y)=(X-Y)/(\log X-\log Y)$ and $XY/\ell(X,Y)$,
with the continuous value $\ell(X,X)=X$.
The identities that implement these operations are
\begin{equation}\label{int:log-means}
 \frac1{\ell(X,Y)}=\int_0^1\frac{dt}{(1-t)X+tY},
 \qquad
 \frac{XY}{\ell(X,Y)}=\int_0^1\frac{dt}{t/X+(1-t)/Y}.
\end{equation}
For the first identity, a positive quadrature followed by inversion
is a finite parallel sum: if $q_j>0$ and $V_j>0$, then
\[
 \left(\sum_{j=1}^M\frac{q_j}{V_j}\right)^{-1}
 = (V_1/q_1):\cdots:(V_M/q_M).
\]
Here $V_j=(1-t_j)X+t_jY$ already satisfies comparison. The second
identity is a positive average of weighted parallel sums.  Pointwise limits give the integrals; zero coefficients
are omitted or obtained by a limit.  Thus no rule asserting preservation
under multiplication of weights is being used.

Apply these two operations to each opposite-exponent pair in
\eqref{int:four-weights}.  Since their logarithmic mean is $m$, we obtain
comparison for the three families
\begin{equation}\label{int:basic-weights}
 m,\qquad ms,\qquad mt,
 \qquad s=f(x)^2,\quad t=f(y)^2.
\end{equation}
The next step constructs one further weight from these three.

\subsection{A Stieltjes representation}

The function $x\mapsto f(x)^2$ decreases from $1$ to $0$ on $[0,\infty)$.
Define $B$ on $(0,1]$ by
\begin{equation}\label{int:B-definition}
 B\bigl(f(x)^2\bigr)=b(x).
\end{equation}
The scalar quadratic \eqref{int:quadratic}, used at the end of the proof,
determines the candidate for the additional weight. It pairs the
four original weights on $\xi\mp\phi$ and $\xi\pm\phi$ and adds
$E|\phi|^2$. For a candidate $E=me>0$, write $a=b(x)+b(y)$. Its minimum is
\[
 \frac m2\left(a-\frac{(x-y)^2}{a+2e}\right)|\xi|^2.
\]
For $x^2\ne y^2$, the addition formula gives
$F/m=(yb(x)+xb(y))/(x+y)$. Equating the minimum with $F|\xi|^2$ therefore
requires, within this quadratic construction,
\[
 (a+2e)(b(x)-b(y))=x^2-y^2,
 \qquad
 E=\frac{m(t-s)}{2(B(s)-B(t))},
 \quad s=f(x)^2,\quad t=f(y)^2,
\]
where the last equality uses $b(z)^2=f(z)^2+z^2$. The final minimization
will verify the value in all cases. It remains to obtain comparison for
this candidate from the known weights. Its reciprocal has the integral
form, initially for $0<s,t<1$,
\begin{equation}\label{int:E-integral}
 \frac1E=\frac2m\int_0^1
       -B'\bigl(\theta s+(1-\theta)t\bigr)\,d\theta,
 \qquad s=f(x)^2,\quad t=f(y)^2.
\end{equation}
The lemma below will show that this defines a positive continuous weight
also when $s=1$ or $t=1$. A Stieltjes representation of $-B'$ turns
the reciprocal into a limit of positive combinations of $1/m$ and
reciprocals of
$\theta(ms)+(1-\theta)(mt)+rm$, with $r\ge0$.
The parallel-sum argument above will then construct $E$ from the three
known weights. We now prove the required representation.

\begin{lemma}\label{int:stieltjes}
The function $B$ extends holomorphically to the upper half-plane and
across the positive real axis, is real and strictly decreasing there,
and satisfies $B(1)=1$, $B'(1)=-1$.  There are a constant $c\ge0$ and a
positive measure $\nu$ on $[0,\infty)$ such that
\begin{equation}\label{int:stieltjes-form}
 -B'(s)=c+\int_{[0,\infty)}\frac{d\nu(r)}{s+r}
 \quad(s>0),\qquad
 \int_{[0,\infty)}\frac{d\nu(r)}{1+r}<\infty.
\end{equation}
In particular $-B'(s)>0$ for every $s>0$.
\end{lemma}

\begin{proof}
For $x>0$, put $s=f(x)^2\in(0,1)$. Differentiating $s=f(x)^2$ and
$B(s)=b(x)$ gives
\[
 \frac{ds}{dx}=\frac{2s(1-B)}{x},\qquad
 \frac{db}{dx}=\frac{B-s}{x}.
\]
Here $B(s)=x\coth x>1$, so we may set
\[
 R(s)=\frac{1-s}{s(B(s)-1)}.
\]
The derivatives give
\begin{equation}\label{int:B-derivative}
 -B'(s)=\frac{B(s)-s}{2s(B(s)-1)}
       =\frac12\left(\frac1s+R(s)\right).
\end{equation}
The term $1/s$ is already a Stieltjes kernel, so it remains to obtain
such a representation for $R$. We first construct the extension of $B$
using
\begin{equation}\label{int:u-parametrization}
 s=\left(\frac{u}{\sin u}\right)^2,\qquad B(s)=u\cot u.
\end{equation}
For $q>0$ let $p_*(q)\in(\pi/2,\pi)$ be the unique solution of
\[
 p_*(q)\tan p_*(q)=-q\tanh q,
\]
and put
\[
 D=\{p+iq:q>0,\ 0<p<p_*(q)\},\qquad
 \mathcal Q=\{z:\operatorname{Re}z>0,\ \operatorname{Im}z<0\}.
\]
Existence, uniqueness and continuity of $p_*$ follow because
$p\tan p$ increases strictly from $-\infty$ to $0$ on $(\pi/2,\pi)$:
its derivative is $(p+\sin p\cos p)/\cos^2p>0$.
Thus $D$ is connected.  The map $a(u)=\sin u/u$ takes $D$ into
$\mathcal Q$, as follows from
\begin{align*}
 \operatorname{Re}a(p+iq)
 &=\frac{p\sin p\cosh q+q\cos p\sinh q}{p^2+q^2}>0,\\
 \operatorname{Im}a(p+iq)
 &=\frac{p\cos p\sinh q-q\sin p\cosh q}{p^2+q^2}<0.
\end{align*}
The first sign is precisely the boundary equation when $p>\pi/2$,
and is immediate for $p\le\pi/2$.  For the second, when $p<\pi/2$ use
$\tan p/p>1>\tanh q/q$; for $p\ge\pi/2$ the sign is immediate.

This map is proper into $\mathcal Q$.  Its finite boundary pieces map
to the coordinate axes: $p=0$ and $q=0$ map to the positive real axis,
and $p=p_*(q)$ to the negative imaginary axis.  The limiting corners
$u=0,\pi$ map to $1,0$, respectively.  Moreover,
\[
 |a(p+iq)|\ge\frac{\sinh q}{\sqrt{\pi^2+q^2}}\longrightarrow\infty
 \qquad(q\longrightarrow\infty)
\]
uniformly for $0\le p\le\pi$.  Thus the inverse image of a compact
subset of $\mathcal Q$ cannot approach the boundary or infinity.
There are no critical points in $D$, because
\begin{equation}\label{int:anti-pick-sign}
 \operatorname{Im}(u\cot u)
 =\frac{q\sin(2p)-p\sinh(2q)}{\cosh(2q)-\cos(2p)}<0,
\end{equation}
whereas $a'(u)=0$ would imply $u\cot u=1$.
The strict sign follows from $\sin(2p)\le2p$ and $\sinh(2q)>2q$.

For completeness, properness and the inverse function theorem imply
that $a(D)$ is both open and relatively closed in $\mathcal Q$, hence
is all of $\mathcal Q$.  Each fiber is finite, and local inverse
neighborhoods at its points give an evenly covered neighborhood: an
additional inverse point arbitrarily close to the chosen image point
would, by properness, accumulate at that fiber.  Thus $a$ is a covering
map.  A connected covering of the simply connected quadrant has one
sheet: inverse branches continue along paths by successive evenly
covered neighborhoods, and homotopies make the continuations
independent of the path.  Therefore $a:D\to\mathcal Q$ has a
holomorphic inverse.

For $\operatorname{Im}s>0$, take the branch of $s^{-1/2}$ in
$\mathcal Q$ and define $u=a^{-1}(s^{-1/2})$.  Equations
\eqref{int:u-parametrization} and \eqref{int:anti-pick-sign} give a
holomorphic $B$ with $\operatorname{Im}B<0$.
To check its continuation at $s_0>0$, the preceding bound on $|a|$
bounds $q$ as $s\to s_0$ from above.  Every cluster point of $u$ lies
on $p=0$ or $q=0$, since the curved boundary has purely imaginary
image.  On these pieces $\sinh q/q$ increases strictly from $1$ to
$\infty$, and $\sin p/p$ decreases strictly from $1$ to $0$ for
$0<p<\pi$.  These observations identify a unique boundary limit:
$u=iq$ if $0<s_0<1$, and $u=p\in(0,\pi)$ if $s_0>1$.
The derivatives there are nonzero, so local analytic inversion gives
continuation across $s_0$.  At $s_0=1$ use $u^2$ as parameter:
\[
 \left(\frac{u}{\sin u}\right)^2=1+\frac{u^2}{3}+O(u^4),
 \qquad u\cot u=1-\frac{u^2}{3}+O(u^4).
\]
This gives $B(1)=1$, $B'(1)=-1$.  The two real parametrizations
also show that $B$ is strictly decreasing on $(0,\infty)$ and agrees
with \eqref{int:B-definition}.

We next recall the classical Herglotz representation
\cite{Herglotz1911}, with a positive-harmonic proof in the form needed
below.  If $P$ is holomorphic with $\operatorname{Im}P\ge0$ in
the upper half-plane and extends holomorphically with real values
across $(0,\infty)$, then
\begin{equation}\label{int:herglotz}
 P(s)=a_0+\beta s+\int_{(-\infty,0]}
 \left(\frac1{r-s}-\frac{r}{1+r^2}\right)d\mu(r),
 \qquad \beta\ge0,
\end{equation}
where $a_0\in\R$, $\mu\ge0$ and
$\int(1+r^2)^{-1}d\mu(r)<\infty$.
Indeed, under $s=i(1+z)/(1-z)$ its imaginary part becomes a
nonnegative harmonic function $v$ on the disk.  The measures
$v(\varrho e^{i\theta})d\theta/(2\pi)$ have the fixed mass $v(0)$.
A weak limit as $\varrho\uparrow1$, together with the Poisson formula
on each smaller disk, represents $v$ by a positive boundary measure.
The atom at $z=1$ contributes $\beta\operatorname{Im}s$; at the
other boundary points the Poisson kernel becomes
$(1+r^2)\operatorname{Im}s/|r-s|^2$.  Incorporate the factor
$1+r^2$ into $d\mu(r)$.  The measure has no mass on the arcs
corresponding to $r>0$, because $v$ tends uniformly to zero on their
compact subarcs, by the assumed continuation.  The analytic
integral in \eqref{int:herglotz} converges locally uniformly; its
imaginary part is the one just obtained.  The remaining holomorphic
function has zero imaginary part and is a real constant.  This
proves \eqref{int:herglotz}.

Apply \eqref{int:herglotz} to $P=-B$ and subtract the value at $1$.
The function
\begin{equation}\label{int:G}
 G(s)=\frac{B(s)-1}{1-s}
 =\beta+\int_{(-\infty,0]}\frac{d\mu(r)}{(1-r)(s-r)}
\end{equation}
has a removable value $G(1)=1$.  It is positive on the positive real
axis, and $H(s)=sG(s)$ has nonnegative imaginary part in the upper
half-plane: this follows termwise from $\beta\ge0$ and
$\operatorname{Im}(s/(s-r))\ge0$ for $r\le0$.
The function $H$ has no zero there.  Otherwise the nonnegative
harmonic function $\operatorname{Im}H$ would attain zero inside,
and the minimum principle would force $H$ to be a real constant;
its positive real-axis values rule out the constant zero.
Thus
\[
 R(s)=\frac1{sG(s)}
\]
continues the real function $R$ above holomorphically to the upper
half-plane and across $(0,\infty)$,
is positive on that interval, and satisfies
$\operatorname{Im}(-R)\ge0$.

Apply \eqref{int:herglotz} once more, now to $-R$, with coefficients
$\widetilde\beta,\widetilde\mu$.  For real $s>1$ it gives
\begin{align*}
 0\le R(1)-R(s)
 &=\widetilde\beta(s-1)
   +\int_{(-\infty,0]}
      \left(\frac1{1-r}-\frac1{s-r}\right)d\widetilde\mu(r)\\
 &\le R(1).
\end{align*}
Consequently $\widetilde\beta=0$, and monotone convergence gives
$\int(1-r)^{-1}d\widetilde\mu(r)\le R(1)$.  We may therefore write
\begin{equation}\label{int:R-stieltjes}
 R(s)=c_0+\int_{(-\infty,0]}\frac{d\widetilde\mu(r)}{s-r},
 \qquad
 c_0=R(1)-\int_{(-\infty,0]}\frac{d\widetilde\mu(r)}{1-r}\ge0.
\end{equation}
Initially this identity follows on the positive axis; local uniform
convergence and analytic continuation give it in the upper half-plane
as well.

The holomorphic functions $-B'$ and $\tfrac12(1/s+R(s))$ agree on
$0<s<1$ by \eqref{int:B-derivative}. The identity theorem extends their
equality to the upper half-plane and the positive real axis. Equation
\eqref{int:R-stieltjes}, after replacing $r$ by $-r$, now proves
\eqref{int:stieltjes-form}; the term $1/(2s)$ contributes an atom
of mass $1/2$ at zero and ensures strict positivity.
\end{proof}

\subsection{The additional weight and the final minimization}

The lemma extends the same integral definition \eqref{int:E-integral}
continuously and positively to all $s,t>0$, including the values used in
\eqref{int:basic-weights}. The fundamental theorem of calculus gives
\begin{equation}\label{int:E}
 E=
 \begin{cases}
 \displaystyle\frac{m(t-s)}{2(B(s)-B(t))},&s\ne t,\\[6pt]
 \displaystyle-\frac{m}{2B'(s)},&s=t.
\end{cases}
\end{equation}
Lemma~\ref{int:stieltjes} shows that $E$ is obtainable by the
comparison-preserving operations already established.  Indeed, after
substitution of \eqref{int:stieltjes-form}, every reciprocal term is
of the form
\[
 \frac1{m(\theta s+(1-\theta)t+r)}
 =\frac1{\theta(ms)+(1-\theta)(mt)+rm},
\]
and the constant term is a multiple of $1/m$.
Positive finite quadratures followed by inversion are parallel sums
of positive multiples of these positive linear combinations of
$m,ms,mt$.  Truncating the $r$-integral and refining partitions gives
pointwise convergence for every $s,t>0$: on each bounded interval the
integrands are continuous, and the tail is controlled by
$\int(1+r)^{-1}d\nu(r)<\infty$.
The source finite-parameter hypothesis and target Fatou argument
therefore prove comparison for $E$.

We complete the proof of Theorem~\ref{thm:interpolation} by an exact
scalar minimization.  For $\xi,\phi\in\C$ put
\begin{equation}\label{int:quadratic}
 \mathcal Q(\xi,\phi)=
 \frac14\sum_{\pm}mf(x)e^{\pm x}|\xi\mp\phi|^2
 +\frac14\sum_{\pm}mf(y)e^{\pm y}|\xi\pm\phi|^2
 +E|\phi|^2.
\end{equation}
Write $a=b(x)+b(y)$ and $e=E/m$.  Since
$f(x)e^{\pm x}=b(x)\pm x$, expansion gives
\[
 \mathcal Q(\xi,\phi)
 =\frac{ma}{2}|\xi|^2
  +m\left(\frac a2+e\right)|\phi|^2
  -m(x-y)\operatorname{Re}(\overline\xi\phi).
\]
The unique minimizer and its value are
\begin{equation}\label{int:minimizer}
 \phi=\frac{x-y}{a+2e}\xi,
 \qquad
 \min_\phi\mathcal Q(\xi,\phi)
 =\frac m2\left(a-\frac{(x-y)^2}{a+2e}\right)|\xi|^2.
\end{equation}
If $x^2\ne y^2$, the identity $b(z)^2=f(z)^2+z^2$ and
\eqref{int:E} give
\[
 (a+2e)(b(x)-b(y))=x^2-y^2.
\]
Hence the minimum in \eqref{int:minimizer} is
\begin{equation}\label{int:minimum-value}
 m\frac{y b(x)+x b(y)}{x+y}|\xi|^2
 =m\frac{f(x)f(y)}{f(x+y)}|\xi|^2,
\end{equation}
where the second equality is the addition formula for $\sinh$,
with continuous values when $x$ or $y$ is zero.

No singular minimization is hidden when $x^2=y^2$.
If $y=x$, the minimizer is $\phi=0$ and the value is
$mb(x)|\xi|^2=mf(x)^2|\xi|^2/f(2x)$.
If $y=-x\ne0$, write $s=f(x)^2$ and $b=b(x)$.
Equations \eqref{int:E} and \eqref{int:B-derivative} yield
\[
 e=\frac{s(b-1)}{b-s},\qquad
 b+e=\frac{x^2}{b-s}.
\]
Thus \eqref{int:minimizer} gives
$m(b-x^2/(b+e))|\xi|^2=ms|\xi|^2$, as required.
At $x=y=0$, we have $E=m/2$, $\phi=0$, and the value is $m|\xi|^2$.
This proves \eqref{int:minimum-value} in every case.

For a source test $Z$, let $\Phi$ be its coordinatewise minimizer in
\eqref{int:minimizer}.  It belongs to $\mathcal T$ by the invariance
assumption.  Apply the four assumed comparisons to $Z\mp\Phi$ and
$Z\pm\Phi$, and the comparison for $E$ to $\Phi$.
After summing, the source side is $Q_{F,0}(Z)$ by
\eqref{int:minimum-value}, whereas the target side is
\[
 \sum_i w_i\sum_\alpha
   \mathcal Q_i\bigl((L_iZ)_\alpha,(L_i\Phi)_\alpha\bigr)
 \ge\sum_i w_i Q_{F,i}(L_iZ).
\]
This is the claimed comparison, and proves
Theorem~\ref{thm:interpolation}.
